\chapter{Glossário de termos metodológicos}
\label{ap:glossario_metodologico}

Este apêndice reúne definições curtas de termos metodológicos e analíticos recorrentes na dissertação. O objetivo é oferecer apoio de leitura a públicos não especializados em modelagem, simulação e análise quantitativa, preservando no texto principal a fluidez da argumentação.

\begin{description}
    \item[Algoritmo] Sequência ordenada de etapas usada para executar um cálculo, resolver um problema ou gerar resultados a partir de dados de entrada.

    \item[Analítico] Relativo à decomposição de um problema em partes e relações formais, de modo a permitir interpretação estruturada dos resultados.

    \item[Cenário] Combinação específica de hipóteses, parâmetros e condições de análise usada para representar uma situação possível no modelo.

    \item[Cenário-base] Cenário de referência a partir do qual os demais cenários são comparados. Nesta dissertação, ele corresponde ao uso do \emph{bunker fuel} convencional nas condições operacionais da rota analisada.

    \item[Critério de viabilidade] Regra adotada para classificar um cenário como viável ou não. No presente estudo, esse critério combina simultaneamente um requisito mínimo de redução de emissões e um requisito econômico-regulatório.

    \item[Custo marginal de abatimento] Indicador que relaciona o custo adicional de uma alternativa à quantidade de emissões que ela evita. Ele mostra quanto custa reduzir uma tonelada adicional de dióxido de carbono equivalente.

    \item[Determinístico] Tipo de modelagem em que os resultados decorrem diretamente dos valores de entrada definidos, sem incorporar sorteios aleatórios ou distribuições probabilísticas.

    \item[Domínio do modelo] Conjunto de combinações admissíveis entre variáveis e parâmetros que o modelo efetivamente examina.

    \item[Hipótese operacional] Formulação testável, expressa nos termos do próprio modelo, que orienta a verificação principal do estudo.

    \item[Modelo de simulação] Representação simplificada de um sistema real usada para calcular o comportamento esperado desse sistema sob diferentes condições.

    \item[Parâmetro] Quantidade tratada como dada no modelo, usada para caracterizar condições de operação, preços, fatores de emissão ou demais propriedades relevantes.

    \item[Robustez] Grau em que um resultado se mantém estável quando certas hipóteses ou parâmetros do modelo são alterados.

    \item[Sensibilidade] Medida de quanto os resultados do modelo mudam quando uma ou mais entradas são alteradas.

    \item[Simulação] Procedimento de execução do modelo para reproduzir, de forma simplificada, o comportamento de um sistema em condições específicas.

    \item[Estocástico] Tipo de modelagem que incorpora incerteza explicitamente, representando variáveis por distribuições probabilísticas ou mecanismos aleatórios.

    \item[Variável] Grandeza cujo valor pode mudar entre cenários ou ao longo da modelagem, influenciando ou representando os resultados do sistema analisado.

    \item[Varredura paramétrica] Procedimento em que um ou mais parâmetros são alterados sistematicamente dentro de um intervalo predefinido para observar como os resultados respondem a essas mudanças.
\end{description}
